\section{Discussion}
\label{sec:discussion}

\subsection{Interpretation of the hybrid correction}

The principal positive result is internally consistent: the residual network
reduced continuous depth and velocity errors for one-step predictions, and a
relaxed autoregressive coupling retained a smaller but positive depth-RMSE
improvement over the pure JAX trajectory.  The reduction from 66.9010\%
one-step depth skill to 17.7405\% on the held-out rollout is expected evidence
of error accumulation and distribution shift when corrected states are fed
back into the physics solver.  It also justifies the conservative correction
relaxation of 0.05 rather than direct unit-strength feedback.

The depth-volume projection is a useful physical safeguard.  It prevents the
network from improving RMSE merely by adding or removing net water after each
physics interval.  Nevertheless, volume preservation alone does not enforce
local momentum conservation, energy behaviour, or consistency with the
shallow-water fluxes.  The 12.0528~m\(^3\) maximum hybrid--JAX volume
difference is small relative to JAX volume, while both structured-grid
trajectories retain a substantial gridded volume difference from ANUGA.  That
difference can arise from discretisation, raster masking, forcing
representation, or solver numerics and should be investigated before treating
the models as interchangeable.

\subsection{Continuous error versus flood extent}

The one-step network improved RMSE and MAE but slightly reduced CSI at the
0.01~m flood threshold.  These metrics answer different questions.  A
continuous loss can favour broad reductions in small depth error, whereas CSI
depends only on threshold crossings.  Here the corrected model removed many
false positives but also removed most true positives.  Therefore, the current
model cannot be claimed to improve binary flood-extent forecasting at the
configured threshold.  Future training could include a differentiable
extent-sensitive term, but its threshold and weight would be a new modelling
choice requiring validation on deeper and independent events.

\subsection{Hydraulic and data limitations}

The terrain is synthetic and has no surveyed vertical datum.  LULC classes
stand in for unresolved roads, buildings, and resistance, and the preserved
waterlogging depressions are seeded scenarios rather than observed locations.
No gauge levels, discharge measurements, observed inundation maps, or
independent storm events are present.  ANUGA is consequently a numerical
reference, not ground truth.

The finalized reflective-boundary choice gives the three model branches a
common comparison basis, but closes every perimeter route.  If
\texttt{2Dboundary\_1} or another segment was intended as an outlet, the
experiment underestimates outflow and can over-retain rainfall.  The one-hour
window compounds this limitation: it covers only 2.9412\% of the rainfall
record and ends before the recorded peak.  Maximum depths are only centimetric,
which makes the 0.01~m extent score particularly sensitive.  Results therefore
characterise an early-event, closed-domain pipeline experiment, not
operational Gurugram flood skill or the complete stated return-period event.

\subsection{Computational claims}

JAX execution and residual training completed on CPU because no usable NVIDIA
driver was available.  Although the solver and coupled loop are JIT- and
scan-compatible, code structure alone is not evidence of GPU acceleration.
The matched benchmark includes appropriate synchronisation, repeated
measurements, separate compile cost, and peak-memory accounting, but its
publication configuration has not run.  Claims of real-time performance or a
JAX speed advantage must therefore wait for the same-GPU experiment.
